% GENERATED FILE. Edit canonical lessons, then run npm run generate:books.
% canonical-source-hash: fnv1a64:b2540d3f1e81fa19
% canonical-lessons: HI-C338-pari, HI-C338-jadu, HI-C338-bhut, HI-C338-dayan, HI-C338-camatkar

\chapter{Fairies and Ghosts}
\label{ch:hi-fairies-and-ghosts}
\hlchaptermodality{338}

\begin{chapteropening}
\textbf{By the end of this chapter:} \emph{I can tell a story with magic in it.}

Complete the last lesson of chapter 338: I can tell a story with magic in it.
\end{chapteropening}

\section[\texorpdfstring{\hi{परी} (parī)}{परी (parī)}]{\hi{परी} (parī) — a fairy}

\label{lesson:HI-C338-pari}



\begin{quote}
Before the new one: say the Hindi for a jewel, then the Hindi for to bite (of a snake).
\end{quote}

\subsection*{You'll want to know: \hi{परी}}
\textbf{\hi{परी}} — \emph{parī} — \textquotedblleft{}a fairy\textquotedblright{}.

\begin{quote}
\hi{कहानी} \hi{में} \hi{एक} \hi{परी} \hi{आसमान} \hi{से} \hi{आई।}
\end{quote}

\begin{quote}
\emph{kahānī meṁ ek parī āsmān se āī.}
\end{quote}

\begin{quote}
In the story a fairy came down from the sky.
\end{quote}

Feminine.

\begin{cousinweb}
is usually traced to Persian \emph{parī}, \textquotedblleft{}a fairy\textquotedblright{}, a winged spirit of Persian tales. A very beautiful girl is also called a \textbf{\hi{परी}}.
\end{cousinweb}

\begin{grammarlens}[title={What you've built}]
One more word for telling what happened.
\end{grammarlens}

\subsection*{Your turn}
\begin{itemize}
\raggedright
  \item \emph{Say it:} \emph{parī}
  \item \emph{Say it:} \emph{parī}, once more
  \item \emph{Say it:} \emph{kahānī meṁ ek parī āsmān se āī.} — In the story a fairy came down from the sky.
  \item \emph{Recall:} say the Hindi for a jewel, then the Hindi for to bite (of a snake), then say \emph{parī} again
\end{itemize}

\begin{tcolorbox}[breakable,colback=teal!4,colframe=teal!35!black,title={Before you move on}]
What does \hi{परी} mean? (\textquotedblleft{}A fairy\textquotedblright{}.) Say it once more.
\end{tcolorbox}
\section[\texorpdfstring{\hi{जादू} (jādū)}{जादू (jādū)}]{\hi{जादू} (jādū) — magic}

\label{lesson:HI-C338-jadu}



\begin{quote}
Before the new one: say the Hindi for to bite (of a snake), then the Hindi for a fairy.
\end{quote}

\subsection*{You'll want to know: \hi{जादू}}
\textbf{\hi{जादू}} — \emph{jādū} — \textquotedblleft{}magic\textquotedblright{}.

\begin{quote}
\hi{परी} \hi{ने} \hi{जादू} \hi{से} \hi{पत्थर} \hi{को} \hi{फूल} \hi{बनाया।}
\end{quote}

\begin{quote}
\emph{parī ne jādū se patthar ko phūl banāyā.}
\end{quote}

\begin{quote}
The fairy used magic to turn a stone into a flower.
\end{quote}

Masculine.

\begin{cousinweb}
\textbf{\hi{जादू}} (\emph{jādū}) is usually traced to Persian \emph{jādū}, \textquotedblleft{}magic\textquotedblright{}. \textbf{\hi{जादू} \hi{से}} is by magic.
\end{cousinweb}

\begin{grammarlens}[title={What you've built}]
One more word for telling what happened.
\end{grammarlens}

\subsection*{Your turn}
\begin{itemize}
\raggedright
  \item \emph{Say it:} \emph{jādū}
  \item \emph{Say it:} \emph{jādū}, once more
  \item \emph{Say it:} \emph{parī ne jādū se patthar ko phūl banāyā.} — The fairy used magic to turn a stone into a flower.
  \item \emph{Recall:} say the Hindi for to bite (of a snake), then the Hindi for a fairy, then say \emph{jādū} again
\end{itemize}

\begin{tcolorbox}[breakable,colback=teal!4,colframe=teal!35!black,title={Before you move on}]
What does \hi{जादू} mean? (\textquotedblleft{}Magic\textquotedblright{}.) Say it once more.
\end{tcolorbox}
\section[\texorpdfstring{\hi{भूत} (bhūt)}{भूत (bhūt)}]{\hi{भूत} (bhūt) — a ghost}

\label{lesson:HI-C338-bhut}



\begin{quote}
Before the new one: say the Hindi for a fairy, then the Hindi for magic.
\end{quote}

\subsection*{You'll want to know: \hi{भूत}}
\textbf{\hi{भूत}} — \emph{bhūt} — \textquotedblleft{}a ghost\textquotedblright{}.

\begin{quote}
\hi{लोग} \hi{कहते} \hi{थे} \hi{कि} \hi{उस} \hi{पुराने} \hi{महल} \hi{में} \hi{भूत} \hi{रहता} \hi{है।}
\end{quote}

\begin{quote}
\emph{log kahte the ki us purāne mahal meṁ bhūt rahtā hai.}
\end{quote}

\begin{quote}
People used to say that a ghost lived in that old palace.
\end{quote}

Masculine.

\begin{cousinweb}
\textbf{\hi{भूत}} (\emph{bhūt}) is Sanskrit \emph{bhūta}, \textquotedblleft{}what has been\textquotedblright{}, and so a being and a spirit (Monier-Williams). The same word means the past in grammar; in stories it is a ghost.
\end{cousinweb}

\begin{grammarlens}[title={What you've built}]
One more word for telling what happened.
\end{grammarlens}

\subsection*{Your turn}
\begin{itemize}
\raggedright
  \item \emph{Say it:} \emph{bhūt}
  \item \emph{Say it:} \emph{bhūt}, once more
  \item \emph{Say it:} \emph{log kahte the ki us purāne mahal meṁ bhūt rahtā hai.} — People used to say that a ghost lived in that old palace.
  \item \emph{Recall:} say the Hindi for a fairy, then the Hindi for magic, then say \emph{bhūt} again
\end{itemize}

\begin{tcolorbox}[breakable,colback=teal!4,colframe=teal!35!black,title={Before you move on}]
What does \hi{भूत} mean? (\textquotedblleft{}A ghost\textquotedblright{}.) Say it once more.
\end{tcolorbox}
\section[\texorpdfstring{\hi{डायन} (ḍāyan)}{डायन (ḍāyan)}]{\hi{डायन} (ḍāyan) — a witch}

\label{lesson:HI-C338-dayan}



\begin{quote}
Before the new one: say the Hindi for magic, then the Hindi for a ghost.
\end{quote}

\subsection*{You'll want to know: \hi{डायन}}
\textbf{\hi{डायन}} — \emph{ḍāyan} — \textquotedblleft{}a witch\textquotedblright{}.

\begin{quote}
\hi{कहानी} \hi{की} \hi{डायन} \hi{जंगल} \hi{में} \hi{अकेली} \hi{रहती} \hi{थी।}
\end{quote}

\begin{quote}
\emph{kahānī kī ḍāyan jaṅgal meṁ akelī rahtī thī.}
\end{quote}

\begin{quote}
The witch in the story lived alone in the forest.
\end{quote}

Feminine.

\begin{cousinweb}
\textbf{\hi{डायन}} (\emph{ḍāyan}) goes back to Sanskrit \emph{ḍākinī}, a female attendant of the goddess Kālī (CDIAL 5542), where CDIAL lists Hindi \emph{ḍāin} and \emph{ḍākin}: in village tales, a witch.
\end{cousinweb}

\begin{grammarlens}[title={What you've built}]
One more word for telling what happened.
\end{grammarlens}

\subsection*{Your turn}
\begin{itemize}
\raggedright
  \item \emph{Say it:} \emph{ḍāyan}
  \item \emph{Say it:} \emph{ḍāyan}, once more
  \item \emph{Say it:} \emph{kahānī kī ḍāyan jaṅgal meṁ akelī rahtī thī.} — The witch in the story lived alone in the forest.
  \item \emph{Recall:} say the Hindi for magic, then the Hindi for a ghost, then say \emph{ḍāyan} again
\end{itemize}

\begin{tcolorbox}[breakable,colback=teal!4,colframe=teal!35!black,title={Before you move on}]
What does \hi{डायन} mean? (\textquotedblleft{}A witch\textquotedblright{}.) Say it once more.
\end{tcolorbox}
\section[\texorpdfstring{\hi{चमत्कार} (camatkār)}{चमत्कार (camatkār)}]{\hi{चमत्कार} (camatkār) — a miracle; a wonder}

\label{lesson:HI-C338-camatkar}



\begin{quote}
Before the new one: say the Hindi for a ghost, then the Hindi for a witch.
\end{quote}

\subsection*{You'll want to know: \hi{चमत्कार}}
\textbf{\hi{चमत्कार}} — \emph{camatkār} — \textquotedblleft{}a miracle; a wonder\textquotedblright{}.

\begin{quote}
\hi{बच्चे} \hi{का} \hi{बचना} \hi{एक} \hi{चमत्कार} \hi{था।}
\end{quote}

\begin{quote}
\emph{bacce kā bacnā ek camatkār thā.}
\end{quote}

\begin{quote}
The child's survival was a miracle.
\end{quote}

Masculine.

\begin{cousinweb}
\textbf{\hi{चमत्कार}} (\emph{camatkār}) is Sanskrit \emph{camatkāra}, \textquotedblleft{}astonishment, a marvel\textquotedblright{} (Monier-Williams): something that cannot be explained.
\end{cousinweb}

\begin{grammarlens}[title={What you've built}]
One more word for telling what happened.
\end{grammarlens}

\subsection*{Your turn}
\begin{itemize}
\raggedright
  \item \emph{Say it:} \emph{camatkār}
  \item \emph{Say it:} \emph{camatkār}, once more
  \item \emph{Say it:} \emph{bacce kā bacnā ek camatkār thā.} — The child's survival was a miracle.
  \item \emph{Recall:} say the Hindi for a ghost, then the Hindi for a witch, then say \emph{camatkār} again
\end{itemize}

\begin{tcolorbox}[breakable,colback=teal!4,colframe=teal!35!black,title={Before you move on}]
What does \hi{चमत्कार} mean? (\textquotedblleft{}A miracle; a wonder\textquotedblright{}.) Say it once more.
\end{tcolorbox}
